\subsection{Summary of the Results}\label{sec:results-summary}
The TTF return series shows heavy tails, a positive skew, and volatility clustering, and the test split carries these properties through the 2026 disruption. Among the configurations of the ForGAN baseline, a longer training schedule and a larger discriminator raise the excess kurtosis of the generated returns, while the generator hidden size, the noise dimension, a longer condition window, and the input scaling produce no separable improvement. Configuration 8 comes closest to the test split among the stable configurations. None of the nine Fin-GAN branches is separable from the ForGAN baseline on the excess kurtosis. The four branches containing $\mathrm{SR}^{*}$ turn the skewness negative. The $\mathrm{PnL}^{*}_{a}$ and $\mathrm{PnL}^{*}_{a}$ \& $\mathrm{MSE}$ branches come closest to the test split, but their seed spreads are wider than the spread of the ForGAN baseline. Configuration 8 with the ForGAN baseline cost function was selected for pricing.

The selected model reproduces the stylized facts only in part. Its option prices deviate further from the market than the prices of both benchmark models, at an implied volatility error of \msa{0.256}{0.043} against \num{0.119} for the block bootstrap, \num{0.142} for the $\mathrm{iid}$ bootstrap, and \num{0.149} for GBM. The deviation is a systematic underpricing in every moneyness range. The simulated spread of the selected model responds to the returns preceding the valuation date, but it stays below the market on most valuation dates. These results are discussed in Chapter~\ref{sec:discussion}.